%insertfigure(ファイルの名前, 図タイトル, width)
\newcommand{\insertfigure}[3]{
\begin{figure}[htbp]
    \begin{center}
        \includegraphics[width={#3}]{#1}
        \caption[#2]{#2．}
        \label{fig:#2}
    \end{center}
\end{figure}
}

%insertfigure(ファイルの名前, 図タイトル, width，ref)
\newcommand{\insertreffigure}[4]{
\begin{figure}[htbp]
    \begin{center}
        \includegraphics[width={#3}]{#1}
        \caption[#2]{#2．#4より引用．}
        \label{fig:#2}
    \end{center}
\end{figure}
}

%inserttable(表タイトル, 表本体)
\newcommand{\inserttable}[2]{
\begin{table}[htbp]
    \begin{center}
        \caption{#1}
        \label{tb:#1}
        #2
    \end{center}
\end{table}
}

%inserteqnarray(式タイトル, 式本体)
\newcommand{\inserteqnarray}[2]{
\begin{eqnarray}
    \label{eqn:#1}
    #2
\end{eqnarray}
}
